% --- Archon physics formalization source begin ---
% archon:physics
% archon:covers PhyXMiniProblems/problem_phyx_mini_0516.lean
% archon:source-report reports/phyx_mini/problem_phyx_mini_0516.source.json

\chapter{Physics problem phyx\_mini\_0516}
\label{ch:PhyXMiniProblems_problem_phyx_mini_0516}

\paragraph{Problem source.}
\#\# Physical scenario

When a certain photoelectric surface is illuminated with light of different wavelengths, the following stopping potentials are observed as shown in figure:Plot the stopping potential on the vertical axis against the frequency of the light on the horizontal axis.

\#\# Question

Determine the value of Planck's constant \textbackslash{}( h \textbackslash{}) (assuming that the value of \textbackslash{}( e \textbackslash{}) is known).

\#\# Answer choices

- A: A: \textbackslash{}( 6.71\textbackslash{}times 10\textasciicircum{}\{-34\} \textbackslash{}text\{ J\} \textbackslash{}cdot \textbackslash{}text\{s\} \textbackslash{})

- B: B: \textbackslash{}( 6.50\textbackslash{}times 10\textasciicircum{}\{-34\} \textbackslash{}text\{ J\} \textbackslash{}cdot \textbackslash{}text\{s\} \textbackslash{})

- C: C: \textbackslash{}( 6.63\textbackslash{}times 10\textasciicircum{}\{-34\} \textbackslash{}text\{ J\} \textbackslash{}cdot \textbackslash{}text\{s\} \textbackslash{})

- D: D: \textbackslash{}( 6.58\textbackslash{}times 10\textasciicircum{}\{-34\} \textbackslash{}text\{ J\} \textbackslash{}cdot \textbackslash{}text\{s\} \textbackslash{})

\#\# Figure caption (auxiliary; use the image as primary evidence)

The image is a table with two columns. The first column is labeled "Wavelength (nm)," and the second column is labeled "Stopping potential (V)." 

The rows contain the following data pairs:

1. Wavelength: 366 nm, Stopping potential: 1.48 V
2. Wavelength: 405 nm, Stopping potential: 1.15 V
3. Wavelength: 436 nm, Stopping potential: 0.93 V
4. Wavelength: 492 nm, Stopping potential: 0.62 V
5. Wavelength: 546 nm, Stopping potential: 0.36 V
6. Wavelength: 579 nm, Stopping potential: 0.24 V

The table has a simple and clean layout with horizontal lines separating the header and the data rows.

\#\# Dataset metadata

Quantum Phenomena; Numerical Reasoning, Physical Model Grounding Reasoning

\paragraph{Recorded answer/context.}
D: D: \textbackslash{}( 6.58\textbackslash{}times 10\textasciicircum{}\{-34\} \textbackslash{}text\{ J\} \textbackslash{}cdot \textbackslash{}text\{s\} \textbackslash{})

\paragraph{Figure/image path.}
/root/proposal\_for\_physic/science-mango/phyx\_mini\_run/phyx\_data/test\_image/516.png

\paragraph{Formalization target.}
create a compiling Lean file with sorry bodies at `PhyXMiniProblems/problem\_phyx\_mini\_0516.lean`. The Lean declarations must preserve the physical quantities, dimensions or dimensional roles, figure labels, governing-law hypotheses, and final relation expressed by this problem.
Use Mathlib/Physlib names found through LeanExplore where available. If a physics API is missing, introduce faithful local abstractions rather than scalar placeholder aliases.

\begin{theorem}[Physics formalization target]
\label{thm:physics:phyx_mini_0516:target}
\lean{PhyXMiniProblems.ProblemPhyXMini0516.problem_phyx_mini_0516}
\uses{def:physics:phyx-mini-0516:phyxminiproblems-problemphyxmini0516-planckconstantinjouleseconds, def:physics:phyx-mini-0516:phyxminiproblems-problemphyxmini0516-photoelectricexperiment, def:physics:phyx-mini-0516:phyxminiproblems-problemphyxmini0516-matchesphotoelectricscenario, def:physics:phyx-mini-0516:phyxminiproblems-problemphyxmini0516-matchessuppliedstoppingpotentialtable, def:physics:phyx-mini-0516:phyxminiproblems-problemphyxmini0516-matchesrequestedfrequencyplot, def:physics:phyx-mini-0516:phyxminiproblems-problemphyxmini0516-usesroundedreferencevalues, def:physics:phyx-mini-0516:phyxminiproblems-problemphyxmini0516-hasphysicalphotoelectricparameters, def:physics:phyx-mini-0516:phyxminiproblems-problemphyxmini0516-satisfiesphotoelectricgraphlaws, def:physics:phyx-mini-0516:phyxminiproblems-problemphyxmini0516-agreesatdisplayedprecision, def:physics:phyx-mini-0516:phyxminiproblems-problemphyxmini0516-matchesanswerchoice, def:physics:phyx-mini-0516:phyxminiproblems-problemphyxmini0516-recordeddatasetanswer}
The assigned autoformalize agent should translate this physics problem into Lean declarations in the covered file, with theorem and lemma proof bodies written as `by sorry`.
\end{theorem}
\begin{proof}
This is an autoformalization task, not a proof task. Produce statements that can later be proved by the physics prover without weakening the physical model.
\end{proof}

% --- Archon declaration topology begin ---
\section{Declaration topology}
\begin{definition}[energy Dimension]
  \label{def:physics:phyx-mini-0516:phyxminiproblems-problemphyxmini0516-energydimension}
  \lean{PhyXMiniProblems.ProblemPhyXMini0516.energyDimension}
  The MLTQ dimension of energy
\end{definition}
\begin{proof}
  This is the declared model definition of energy Dimension.
\end{proof}
\begin{definition}[electric Potential Dimension]
  \label{def:physics:phyx-mini-0516:phyxminiproblems-problemphyxmini0516-electricpotentialdimension}
  \lean{PhyXMiniProblems.ProblemPhyXMini0516.electricPotentialDimension}
  \uses{def:physics:phyx-mini-0516:phyxminiproblems-problemphyxmini0516-energydimension}
  Electric potential has dimension energy per charge
\end{definition}
\begin{proof}
  This is the declared model definition of electric Potential Dimension.
\end{proof}
\begin{definition}[action Dimension]
  \label{def:physics:phyx-mini-0516:phyxminiproblems-problemphyxmini0516-actiondimension}
  \lean{PhyXMiniProblems.ProblemPhyXMini0516.actionDimension}
  \uses{def:physics:phyx-mini-0516:phyxminiproblems-problemphyxmini0516-energydimension}
  Planck's constant has the dimension of action, energy times time
\end{definition}
\begin{proof}
  This is the declared model definition of action Dimension.
\end{proof}
\begin{definition}[Wavelength Quantity]
  \label{def:physics:phyx-mini-0516:phyxminiproblems-problemphyxmini0516-wavelengthquantity}
  \lean{PhyXMiniProblems.ProblemPhyXMini0516.WavelengthQuantity}
  A nonnegative, unit-independent physical wavelength
\end{definition}
\begin{proof}
  This is the declared model definition of Wavelength Quantity.
\end{proof}
\begin{definition}[Frequency Quantity]
  \label{def:physics:phyx-mini-0516:phyxminiproblems-problemphyxmini0516-frequencyquantity}
  \lean{PhyXMiniProblems.ProblemPhyXMini0516.FrequencyQuantity}
  A nonnegative, unit-independent physical frequency
\end{definition}
\begin{proof}
  This is the declared model definition of Frequency Quantity.
\end{proof}
\begin{definition}[Electric Potential Quantity]
  \label{def:physics:phyx-mini-0516:phyxminiproblems-problemphyxmini0516-electricpotentialquantity}
  \lean{PhyXMiniProblems.ProblemPhyXMini0516.ElectricPotentialQuantity}
  \uses{def:physics:phyx-mini-0516:phyxminiproblems-problemphyxmini0516-electricpotentialdimension}
  A signed, unit-independent electric potential
\end{definition}
\begin{proof}
  This is the declared model definition of Electric Potential Quantity.
\end{proof}
\begin{definition}[Charge Magnitude Quantity]
  \label{def:physics:phyx-mini-0516:phyxminiproblems-problemphyxmini0516-chargemagnitudequantity}
  \lean{PhyXMiniProblems.ProblemPhyXMini0516.ChargeMagnitudeQuantity}
  A nonnegative physical magnitude of electric charge
\end{definition}
\begin{proof}
  This is the declared model definition of Charge Magnitude Quantity.
\end{proof}
\begin{definition}[Energy Quantity]
  \label{def:physics:phyx-mini-0516:phyxminiproblems-problemphyxmini0516-energyquantity}
  \lean{PhyXMiniProblems.ProblemPhyXMini0516.EnergyQuantity}
  \uses{def:physics:phyx-mini-0516:phyxminiproblems-problemphyxmini0516-energydimension}
  A nonnegative physical energy
\end{definition}
\begin{proof}
  This is the declared model definition of Energy Quantity.
\end{proof}
\begin{definition}[Planck Constant Quantity]
  \label{def:physics:phyx-mini-0516:phyxminiproblems-problemphyxmini0516-planckconstantquantity}
  \lean{PhyXMiniProblems.ProblemPhyXMini0516.PlanckConstantQuantity}
  \uses{def:physics:phyx-mini-0516:phyxminiproblems-problemphyxmini0516-actiondimension}
  A nonnegative physical value of Planck's constant
\end{definition}
\begin{proof}
  This is the declared model definition of Planck Constant Quantity.
\end{proof}
\begin{definition}[nonnegative SIReadout]
  \label{def:physics:phyx-mini-0516:phyxminiproblems-problemphyxmini0516-nonnegativesireadout}
  \lean{PhyXMiniProblems.ProblemPhyXMini0516.nonnegativeSIReadout}
  Coherent-SI readout of a nonnegative dimensionful quantity
\end{definition}
\begin{proof}
  This is the declared model definition of nonnegative SIReadout.
\end{proof}
\begin{definition}[signed SIReadout]
  \label{def:physics:phyx-mini-0516:phyxminiproblems-problemphyxmini0516-signedsireadout}
  \lean{PhyXMiniProblems.ProblemPhyXMini0516.signedSIReadout}
  Coherent-SI readout of a signed dimensionful quantity
\end{definition}
\begin{proof}
  This is the declared model definition of signed SIReadout.
\end{proof}
\begin{definition}[wavelength Readout]
  \label{def:physics:phyx-mini-0516:phyxminiproblems-problemphyxmini0516-wavelengthreadout}
  \lean{PhyXMiniProblems.ProblemPhyXMini0516.wavelengthReadout}
  \uses{def:physics:phyx-mini-0516:phyxminiproblems-problemphyxmini0516-wavelengthquantity}
  Read a wavelength in a chosen length unit
\end{definition}
\begin{proof}
  This is the declared model definition of wavelength Readout.
\end{proof}
\begin{definition}[wavelength In Meters]
  \label{def:physics:phyx-mini-0516:phyxminiproblems-problemphyxmini0516-wavelengthinmeters}
  \lean{PhyXMiniProblems.ProblemPhyXMini0516.wavelengthInMeters}
  \uses{def:physics:phyx-mini-0516:phyxminiproblems-problemphyxmini0516-wavelengthquantity, def:physics:phyx-mini-0516:phyxminiproblems-problemphyxmini0516-wavelengthreadout}
  Wavelength readout in metres
\end{definition}
\begin{proof}
  This is the declared model definition of wavelength In Meters.
\end{proof}
\begin{definition}[wavelength In Nanometers]
  \label{def:physics:phyx-mini-0516:phyxminiproblems-problemphyxmini0516-wavelengthinnanometers}
  \lean{PhyXMiniProblems.ProblemPhyXMini0516.wavelengthInNanometers}
  \uses{def:physics:phyx-mini-0516:phyxminiproblems-problemphyxmini0516-wavelengthquantity, def:physics:phyx-mini-0516:phyxminiproblems-problemphyxmini0516-wavelengthreadout}
  Wavelength readout in nanometres, as printed in the table
\end{definition}
\begin{proof}
  This is the declared model definition of wavelength In Nanometers.
\end{proof}
\begin{definition}[frequency In Hertz]
  \label{def:physics:phyx-mini-0516:phyxminiproblems-problemphyxmini0516-frequencyinhertz}
  \lean{PhyXMiniProblems.ProblemPhyXMini0516.frequencyInHertz}
  \uses{def:physics:phyx-mini-0516:phyxminiproblems-problemphyxmini0516-frequencyquantity, def:physics:phyx-mini-0516:phyxminiproblems-problemphyxmini0516-nonnegativesireadout}
  Incident-light frequency readout in hertz
\end{definition}
\begin{proof}
  This is the declared model definition of frequency In Hertz.
\end{proof}
\begin{definition}[stopping Potential In Volts]
  \label{def:physics:phyx-mini-0516:phyxminiproblems-problemphyxmini0516-stoppingpotentialinvolts}
  \lean{PhyXMiniProblems.ProblemPhyXMini0516.stoppingPotentialInVolts}
  \uses{def:physics:phyx-mini-0516:phyxminiproblems-problemphyxmini0516-electricpotentialquantity, def:physics:phyx-mini-0516:phyxminiproblems-problemphyxmini0516-signedsireadout}
  Stopping-potential readout in volts
\end{definition}
\begin{proof}
  This is the declared model definition of stopping Potential In Volts.
\end{proof}
\begin{definition}[speed In Meters Per Second]
  \label{def:physics:phyx-mini-0516:phyxminiproblems-problemphyxmini0516-speedinmeterspersecond}
  \lean{PhyXMiniProblems.ProblemPhyXMini0516.speedInMetersPerSecond}
  Speed readout in metres per second
\end{definition}
\begin{proof}
  This is the declared model definition of speed In Meters Per Second.
\end{proof}
\begin{definition}[charge Magnitude In Coulombs]
  \label{def:physics:phyx-mini-0516:phyxminiproblems-problemphyxmini0516-chargemagnitudeincoulombs}
  \lean{PhyXMiniProblems.ProblemPhyXMini0516.chargeMagnitudeInCoulombs}
  \uses{def:physics:phyx-mini-0516:phyxminiproblems-problemphyxmini0516-chargemagnitudequantity, def:physics:phyx-mini-0516:phyxminiproblems-problemphyxmini0516-nonnegativesireadout}
  Charge-magnitude readout in coulombs
\end{definition}
\begin{proof}
  This is the declared model definition of charge Magnitude In Coulombs.
\end{proof}
\begin{definition}[energy In Joules]
  \label{def:physics:phyx-mini-0516:phyxminiproblems-problemphyxmini0516-energyinjoules}
  \lean{PhyXMiniProblems.ProblemPhyXMini0516.energyInJoules}
  \uses{def:physics:phyx-mini-0516:phyxminiproblems-problemphyxmini0516-energyquantity, def:physics:phyx-mini-0516:phyxminiproblems-problemphyxmini0516-nonnegativesireadout}
  Energy readout in joules
\end{definition}
\begin{proof}
  This is the declared model definition of energy In Joules.
\end{proof}
\begin{definition}[planck Constant In Joule Seconds]
  \label{def:physics:phyx-mini-0516:phyxminiproblems-problemphyxmini0516-planckconstantinjouleseconds}
  \lean{PhyXMiniProblems.ProblemPhyXMini0516.planckConstantInJouleSeconds}
  \uses{def:physics:phyx-mini-0516:phyxminiproblems-problemphyxmini0516-planckconstantquantity, def:physics:phyx-mini-0516:phyxminiproblems-problemphyxmini0516-nonnegativesireadout}
  Planck-constant readout in joule-seconds
\end{definition}
\begin{proof}
  This is the declared model definition of planck Constant In Joule Seconds.
\end{proof}
\begin{definition}[Measurement Row]
  \label{def:physics:phyx-mini-0516:phyxminiproblems-problemphyxmini0516-measurementrow}
  \lean{PhyXMiniProblems.ProblemPhyXMini0516.MeasurementRow}
  The six rows shown in the supplied table, from top to bottom
\end{definition}
\begin{proof}
  This is the declared model definition of Measurement Row.
\end{proof}
\begin{definition}[Table Column]
  \label{def:physics:phyx-mini-0516:phyxminiproblems-problemphyxmini0516-tablecolumn}
  \lean{PhyXMiniProblems.ProblemPhyXMini0516.TableColumn}
  The two columns visible in the primary figure
\end{definition}
\begin{proof}
  This is the declared model definition of Table Column.
\end{proof}
\begin{definition}[Table Column Label]
  \label{def:physics:phyx-mini-0516:phyxminiproblems-problemphyxmini0516-tablecolumnlabel}
  \lean{PhyXMiniProblems.ProblemPhyXMini0516.TableColumnLabel}
  Literal physical roles and unit annotations printed in the table header
\end{definition}
\begin{proof}
  This is the declared model definition of Table Column Label.
\end{proof}
\begin{definition}[Plot Axis]
  \label{def:physics:phyx-mini-0516:phyxminiproblems-problemphyxmini0516-plotaxis}
  \lean{PhyXMiniProblems.ProblemPhyXMini0516.PlotAxis}
  The two axes of the graph requested in the prose
\end{definition}
\begin{proof}
  This is the declared model definition of Plot Axis.
\end{proof}
\begin{definition}[Plot Axis Role]
  \label{def:physics:phyx-mini-0516:phyxminiproblems-problemphyxmini0516-plotaxisrole}
  \lean{PhyXMiniProblems.ProblemPhyXMini0516.PlotAxisRole}
  Physical quantity assigned to each requested plot axis
\end{definition}
\begin{proof}
  This is the declared model definition of Plot Axis Role.
\end{proof}
\begin{definition}[Plot Axis Unit]
  \label{def:physics:phyx-mini-0516:phyxminiproblems-problemphyxmini0516-plotaxisunit}
  \lean{PhyXMiniProblems.ProblemPhyXMini0516.PlotAxisUnit}
  Unit annotation for each requested plot axis
\end{definition}
\begin{proof}
  This is the declared model definition of Plot Axis Unit.
\end{proof}
\begin{definition}[Quantum Process]
  \label{def:physics:phyx-mini-0516:phyxminiproblems-problemphyxmini0516-quantumprocess}
  \lean{PhyXMiniProblems.ProblemPhyXMini0516.QuantumProcess}
  The physical process taking place at the illuminated surface
\end{definition}
\begin{proof}
  This is the declared model definition of Quantum Process.
\end{proof}
\begin{definition}[Stopping Potential Table Figure]
  \label{def:physics:phyx-mini-0516:phyxminiproblems-problemphyxmini0516-stoppingpotentialtablefigure}
  \lean{PhyXMiniProblems.ProblemPhyXMini0516.StoppingPotentialTableFigure}
  \uses{def:physics:phyx-mini-0516:phyxminiproblems-problemphyxmini0516-wavelengthquantity, def:physics:phyx-mini-0516:phyxminiproblems-problemphyxmini0516-electricpotentialquantity, def:physics:phyx-mini-0516:phyxminiproblems-problemphyxmini0516-measurementrow, def:physics:phyx-mini-0516:phyxminiproblems-problemphyxmini0516-tablecolumn, def:physics:phyx-mini-0516:phyxminiproblems-problemphyxmini0516-tablecolumnlabel}
  The primary two-column figure, with dimensionful quantities in each row
\end{definition}
\begin{proof}
  This is the declared model definition of Stopping Potential Table Figure.
\end{proof}
\begin{definition}[Frequency Stopping Potential Plot]
  \label{def:physics:phyx-mini-0516:phyxminiproblems-problemphyxmini0516-frequencystoppingpotentialplot}
  \lean{PhyXMiniProblems.ProblemPhyXMini0516.FrequencyStoppingPotentialPlot}
  \uses{def:physics:phyx-mini-0516:phyxminiproblems-problemphyxmini0516-frequencyquantity, def:physics:phyx-mini-0516:phyxminiproblems-problemphyxmini0516-electricpotentialquantity, def:physics:phyx-mini-0516:phyxminiproblems-problemphyxmini0516-measurementrow, def:physics:phyx-mini-0516:phyxminiproblems-problemphyxmini0516-plotaxis, def:physics:phyx-mini-0516:phyxminiproblems-problemphyxmini0516-plotaxisrole, def:physics:phyx-mini-0516:phyxminiproblems-problemphyxmini0516-plotaxisunit}
  The frequency-versus-stopping-potential graph requested by the problem
\end{definition}
\begin{proof}
  This is the declared model definition of Frequency Stopping Potential Plot.
\end{proof}
\begin{definition}[Photoelectric Experiment]
  \label{def:physics:phyx-mini-0516:phyxminiproblems-problemphyxmini0516-photoelectricexperiment}
  \lean{PhyXMiniProblems.ProblemPhyXMini0516.PhotoelectricExperiment}
  \uses{def:physics:phyx-mini-0516:phyxminiproblems-problemphyxmini0516-frequencyquantity, def:physics:phyx-mini-0516:phyxminiproblems-problemphyxmini0516-chargemagnitudequantity, def:physics:phyx-mini-0516:phyxminiproblems-problemphyxmini0516-energyquantity, def:physics:phyx-mini-0516:phyxminiproblems-problemphyxmini0516-planckconstantquantity, def:physics:phyx-mini-0516:phyxminiproblems-problemphyxmini0516-measurementrow, def:physics:phyx-mini-0516:phyxminiproblems-problemphyxmini0516-quantumprocess, def:physics:phyx-mini-0516:phyxminiproblems-problemphyxmini0516-stoppingpotentialtablefigure, def:physics:phyx-mini-0516:phyxminiproblems-problemphyxmini0516-frequencystoppingpotentialplot}
  Independent quantities in the photoelectric experiment. The Planck constant and work function are fields, not definitions of an answer choice. Their connection to the graph is imposed only by the general physical laws below
\end{definition}
\begin{proof}
  This is the declared model definition of Photoelectric Experiment.
\end{proof}
\begin{definition}[displayed Wavelength In Nanometers]
  \label{def:physics:phyx-mini-0516:phyxminiproblems-problemphyxmini0516-displayedwavelengthinnanometers}
  \lean{PhyXMiniProblems.ProblemPhyXMini0516.displayedWavelengthInNanometers}
  \uses{def:physics:phyx-mini-0516:phyxminiproblems-problemphyxmini0516-measurementrow}
  Nanometre number printed in a given table row
\end{definition}
\begin{proof}
  This is the declared model definition of displayed Wavelength In Nanometers.
\end{proof}
\begin{definition}[displayed Stopping Potential In Volts]
  \label{def:physics:phyx-mini-0516:phyxminiproblems-problemphyxmini0516-displayedstoppingpotentialinvolts}
  \lean{PhyXMiniProblems.ProblemPhyXMini0516.displayedStoppingPotentialInVolts}
  \uses{def:physics:phyx-mini-0516:phyxminiproblems-problemphyxmini0516-measurementrow}
  Volt number printed in a given table row
\end{definition}
\begin{proof}
  This is the declared model definition of displayed Stopping Potential In Volts.
\end{proof}
\begin{definition}[Matches Photoelectric Scenario]
  \label{def:physics:phyx-mini-0516:phyxminiproblems-problemphyxmini0516-matchesphotoelectricscenario}
  \lean{PhyXMiniProblems.ProblemPhyXMini0516.MatchesPhotoelectricScenario}
  \uses{def:physics:phyx-mini-0516:phyxminiproblems-problemphyxmini0516-photoelectricexperiment}
  The prose identifies the experiment as photoelectric emission
\end{definition}
\begin{proof}
  This is the declared model definition of Matches Photoelectric Scenario.
\end{proof}
\begin{definition}[Matches Supplied Stopping Potential Table]
  \label{def:physics:phyx-mini-0516:phyxminiproblems-problemphyxmini0516-matchessuppliedstoppingpotentialtable}
  \lean{PhyXMiniProblems.ProblemPhyXMini0516.MatchesSuppliedStoppingPotentialTable}
  \uses{def:physics:phyx-mini-0516:phyxminiproblems-problemphyxmini0516-wavelengthreadout, def:physics:phyx-mini-0516:phyxminiproblems-problemphyxmini0516-wavelengthinnanometers, def:physics:phyx-mini-0516:phyxminiproblems-problemphyxmini0516-stoppingpotentialinvolts, def:physics:phyx-mini-0516:phyxminiproblems-problemphyxmini0516-stoppingpotentialtablefigure, def:physics:phyx-mini-0516:phyxminiproblems-problemphyxmini0516-displayedwavelengthinnanometers, def:physics:phyx-mini-0516:phyxminiproblems-problemphyxmini0516-displayedstoppingpotentialinvolts}
  Direct transcription of the primary bitmap. These hypotheses give only the two printed headers and the six measured row values; they contain no frequency, fitted slope, Planck constant, or answer choice
\end{definition}
\begin{proof}
  This is the declared model definition of Matches Supplied Stopping Potential Table.
\end{proof}
\begin{definition}[Matches Requested Frequency Plot]
  \label{def:physics:phyx-mini-0516:phyxminiproblems-problemphyxmini0516-matchesrequestedfrequencyplot}
  \lean{PhyXMiniProblems.ProblemPhyXMini0516.MatchesRequestedFrequencyPlot}
  \uses{def:physics:phyx-mini-0516:phyxminiproblems-problemphyxmini0516-photoelectricexperiment}
  The requested graph places the frequency derived for every table row on the horizontal axis and that row's observed stopping potential on the vertical axis. This carries no numerical slope or fitted value
\end{definition}
\begin{proof}
  This is the declared model definition of Matches Requested Frequency Plot.
\end{proof}
\begin{definition}[Uses Rounded Reference Values]
  \label{def:physics:phyx-mini-0516:phyxminiproblems-problemphyxmini0516-usesroundedreferencevalues}
  \lean{PhyXMiniProblems.ProblemPhyXMini0516.UsesRoundedReferenceValues}
  \uses{def:physics:phyx-mini-0516:phyxminiproblems-problemphyxmini0516-speedinmeterspersecond, def:physics:phyx-mini-0516:phyxminiproblems-problemphyxmini0516-chargemagnitudeincoulombs, def:physics:phyx-mini-0516:phyxminiproblems-problemphyxmini0516-photoelectricexperiment}
  The rounded reference values conventionally used in this multiple-choice calculation. They are independent calibrations, not the requested value of Planck's constant. Using Physlib's exact 299792458 m/s and exact elementary charge would change the last displayed digit of this dataset's estimate
\end{definition}
\begin{proof}
  This is the declared model definition of Uses Rounded Reference Values.
\end{proof}
\begin{definition}[Has Physical Photoelectric Parameters]
  \label{def:physics:phyx-mini-0516:phyxminiproblems-problemphyxmini0516-hasphysicalphotoelectricparameters}
  \lean{PhyXMiniProblems.ProblemPhyXMini0516.HasPhysicalPhotoelectricParameters}
  \uses{def:physics:phyx-mini-0516:phyxminiproblems-problemphyxmini0516-wavelengthinmeters, def:physics:phyx-mini-0516:phyxminiproblems-problemphyxmini0516-frequencyinhertz, def:physics:phyx-mini-0516:phyxminiproblems-problemphyxmini0516-speedinmeterspersecond, def:physics:phyx-mini-0516:phyxminiproblems-problemphyxmini0516-chargemagnitudeincoulombs, def:physics:phyx-mini-0516:phyxminiproblems-problemphyxmini0516-energyinjoules, def:physics:phyx-mini-0516:phyxminiproblems-problemphyxmini0516-planckconstantinjouleseconds, def:physics:phyx-mini-0516:phyxminiproblems-problemphyxmini0516-photoelectricexperiment}
  Positivity conditions selecting a physically meaningful experiment
\end{definition}
\begin{proof}
  This is the declared model definition of Has Physical Photoelectric Parameters.
\end{proof}
\begin{definition}[mean Incident Frequency In Hertz]
  \label{def:physics:phyx-mini-0516:phyxminiproblems-problemphyxmini0516-meanincidentfrequencyinhertz}
  \lean{PhyXMiniProblems.ProblemPhyXMini0516.meanIncidentFrequencyInHertz}
  \uses{def:physics:phyx-mini-0516:phyxminiproblems-problemphyxmini0516-frequencyinhertz, def:physics:phyx-mini-0516:phyxminiproblems-problemphyxmini0516-measurementrow, def:physics:phyx-mini-0516:phyxminiproblems-problemphyxmini0516-photoelectricexperiment}
  Arithmetic mean of the six plotted frequencies, in hertz
\end{definition}
\begin{proof}
  This is the declared model definition of mean Incident Frequency In Hertz.
\end{proof}
\begin{definition}[mean Stopping Potential In Volts]
  \label{def:physics:phyx-mini-0516:phyxminiproblems-problemphyxmini0516-meanstoppingpotentialinvolts}
  \lean{PhyXMiniProblems.ProblemPhyXMini0516.meanStoppingPotentialInVolts}
  \uses{def:physics:phyx-mini-0516:phyxminiproblems-problemphyxmini0516-stoppingpotentialinvolts, def:physics:phyx-mini-0516:phyxminiproblems-problemphyxmini0516-measurementrow, def:physics:phyx-mini-0516:phyxminiproblems-problemphyxmini0516-photoelectricexperiment}
  Arithmetic mean of the six measured stopping potentials, in volts
\end{definition}
\begin{proof}
  This is the declared model definition of mean Stopping Potential In Volts.
\end{proof}
\begin{definition}[fitted Stopping Potential Slope]
  \label{def:physics:phyx-mini-0516:phyxminiproblems-problemphyxmini0516-fittedstoppingpotentialslope}
  \lean{PhyXMiniProblems.ProblemPhyXMini0516.fittedStoppingPotentialSlope}
  \uses{def:physics:phyx-mini-0516:phyxminiproblems-problemphyxmini0516-frequencyinhertz, def:physics:phyx-mini-0516:phyxminiproblems-problemphyxmini0516-stoppingpotentialinvolts, def:physics:phyx-mini-0516:phyxminiproblems-problemphyxmini0516-measurementrow, def:physics:phyx-mini-0516:phyxminiproblems-problemphyxmini0516-photoelectricexperiment, def:physics:phyx-mini-0516:phyxminiproblems-problemphyxmini0516-meanincidentfrequencyinhertz, def:physics:phyx-mini-0516:phyxminiproblems-problemphyxmini0516-meanstoppingpotentialinvolts}
  Least-squares slope of stopping potential against frequency. Its scalar unit is volts per hertz, equivalently volt-seconds
\end{definition}
\begin{proof}
  This is the declared model definition of fitted Stopping Potential Slope.
\end{proof}
\begin{definition}[fitted Stopping Potential In Volts]
  \label{def:physics:phyx-mini-0516:phyxminiproblems-problemphyxmini0516-fittedstoppingpotentialinvolts}
  \lean{PhyXMiniProblems.ProblemPhyXMini0516.fittedStoppingPotentialInVolts}
  \uses{def:physics:phyx-mini-0516:phyxminiproblems-problemphyxmini0516-photoelectricexperiment, def:physics:phyx-mini-0516:phyxminiproblems-problemphyxmini0516-meanincidentfrequencyinhertz, def:physics:phyx-mini-0516:phyxminiproblems-problemphyxmini0516-meanstoppingpotentialinvolts, def:physics:phyx-mini-0516:phyxminiproblems-problemphyxmini0516-fittedstoppingpotentialslope}
  The least-squares line evaluated at a frequency readout in hertz
\end{definition}
\begin{proof}
  This is the declared model definition of fitted Stopping Potential In Volts.
\end{proof}
\begin{definition}[Satisfies Photoelectric Graph Laws]
  \label{def:physics:phyx-mini-0516:phyxminiproblems-problemphyxmini0516-satisfiesphotoelectricgraphlaws}
  \lean{PhyXMiniProblems.ProblemPhyXMini0516.SatisfiesPhotoelectricGraphLaws}
  \uses{def:physics:phyx-mini-0516:phyxminiproblems-problemphyxmini0516-wavelengthinmeters, def:physics:phyx-mini-0516:phyxminiproblems-problemphyxmini0516-frequencyinhertz, def:physics:phyx-mini-0516:phyxminiproblems-problemphyxmini0516-speedinmeterspersecond, def:physics:phyx-mini-0516:phyxminiproblems-problemphyxmini0516-chargemagnitudeincoulombs, def:physics:phyx-mini-0516:phyxminiproblems-problemphyxmini0516-energyinjoules, def:physics:phyx-mini-0516:phyxminiproblems-problemphyxmini0516-planckconstantinjouleseconds, def:physics:phyx-mini-0516:phyxminiproblems-problemphyxmini0516-photoelectricexperiment, def:physics:phyx-mini-0516:phyxminiproblems-problemphyxmini0516-fittedstoppingpotentialinvolts}
  The two physical laws used by the intended graph method: every incident wavelength and frequency obey lambda * f = c in vacuum; the least-squares trend line is interpreted by Einstein's photoelectric equation e V\_stop = h f - phi. Neither law contains a numerical slope, a numerical Planck constant, or an answer label
\end{definition}
\begin{proof}
  This is the declared model definition of Satisfies Photoelectric Graph Laws.
\end{proof}
\begin{lemma}[planck constant eq charge mul fitted slope]
  \label{lem:physics:phyx-mini-0516:phyxminiproblems-problemphyxmini0516-planck-constant-eq-charge-mul-fitted-slope}
  \lean{PhyXMiniProblems.ProblemPhyXMini0516.planck_constant_eq_charge_mul_fitted_slope}
  \uses{def:physics:phyx-mini-0516:phyxminiproblems-problemphyxmini0516-chargemagnitudeincoulombs, def:physics:phyx-mini-0516:phyxminiproblems-problemphyxmini0516-planckconstantinjouleseconds, def:physics:phyx-mini-0516:phyxminiproblems-problemphyxmini0516-photoelectricexperiment, def:physics:phyx-mini-0516:phyxminiproblems-problemphyxmini0516-hasphysicalphotoelectricparameters, def:physics:phyx-mini-0516:phyxminiproblems-problemphyxmini0516-fittedstoppingpotentialslope, def:physics:phyx-mini-0516:phyxminiproblems-problemphyxmini0516-satisfiesphotoelectricgraphlaws}
  Einstein's equation identifies the voltage-frequency slope as h / e. This is a general intermediate relation and does not select any displayed number
\end{lemma}
\begin{proof}
  The conclusion follows from the governing relations and figure readouts named in the statement of planck constant eq charge mul fitted slope.
\end{proof}
\begin{definition}[Answer Choice]
  \label{def:physics:phyx-mini-0516:phyxminiproblems-problemphyxmini0516-answerchoice}
  \lean{PhyXMiniProblems.ProblemPhyXMini0516.AnswerChoice}
  Labels of the four answer choices
\end{definition}
\begin{proof}
  This is the declared model definition of Answer Choice.
\end{proof}
\begin{definition}[displayed Planck Constant In Joule Seconds]
  \label{def:physics:phyx-mini-0516:phyxminiproblems-problemphyxmini0516-displayedplanckconstantinjouleseconds}
  \lean{PhyXMiniProblems.ProblemPhyXMini0516.displayedPlanckConstantInJouleSeconds}
  \uses{def:physics:phyx-mini-0516:phyxminiproblems-problemphyxmini0516-answerchoice}
  Planck-constant value printed beside each choice, in joule-seconds
\end{definition}
\begin{proof}
  This is the declared model definition of displayed Planck Constant In Joule Seconds.
\end{proof}
\begin{definition}[answer Half Unit In Joule Seconds]
  \label{def:physics:phyx-mini-0516:phyxminiproblems-problemphyxmini0516-answerhalfunitinjouleseconds}
  \lean{PhyXMiniProblems.ProblemPhyXMini0516.answerHalfUnitInJouleSeconds}
  Half of one unit in the final displayed decimal place of the choices
\end{definition}
\begin{proof}
  This is the declared model definition of answer Half Unit In Joule Seconds.
\end{proof}
\begin{definition}[Agrees At Displayed Precision]
  \label{def:physics:phyx-mini-0516:phyxminiproblems-problemphyxmini0516-agreesatdisplayedprecision}
  \lean{PhyXMiniProblems.ProblemPhyXMini0516.AgreesAtDisplayedPrecision}
  \uses{def:physics:phyx-mini-0516:phyxminiproblems-problemphyxmini0516-answerhalfunitinjouleseconds}
  Agreement at the three-significant-figure precision used by all four choices
\end{definition}
\begin{proof}
  This is the declared model definition of Agrees At Displayed Precision.
\end{proof}
\begin{definition}[Matches Answer Choice]
  \label{def:physics:phyx-mini-0516:phyxminiproblems-problemphyxmini0516-matchesanswerchoice}
  \lean{PhyXMiniProblems.ProblemPhyXMini0516.MatchesAnswerChoice}
  \uses{def:physics:phyx-mini-0516:phyxminiproblems-problemphyxmini0516-planckconstantinjouleseconds, def:physics:phyx-mini-0516:phyxminiproblems-problemphyxmini0516-photoelectricexperiment, def:physics:phyx-mini-0516:phyxminiproblems-problemphyxmini0516-answerchoice, def:physics:phyx-mini-0516:phyxminiproblems-problemphyxmini0516-displayedplanckconstantinjouleseconds, def:physics:phyx-mini-0516:phyxminiproblems-problemphyxmini0516-agreesatdisplayedprecision}
  A modeled Planck constant agrees with a displayed answer choice
\end{definition}
\begin{proof}
  This is the declared model definition of Matches Answer Choice.
\end{proof}
\begin{definition}[recorded Dataset Answer]
  \label{def:physics:phyx-mini-0516:phyxminiproblems-problemphyxmini0516-recordeddatasetanswer}
  \lean{PhyXMiniProblems.ProblemPhyXMini0516.recordedDatasetAnswer}
  \uses{def:physics:phyx-mini-0516:phyxminiproblems-problemphyxmini0516-answerchoice}
  Answer label recorded in the source dataset
\end{definition}
\begin{proof}
  This is the declared model definition of recorded Dataset Answer.
\end{proof}
% --- Archon declaration topology end ---
% --- Archon physics formalization source end ---
